\documentclass[10pt,a4paper]{amsart}
\usepackage[margin=19mm]{geometry}
\usepackage{amsmath,amssymb,mathtools}
\usepackage[scr=rsfs,cal=euler]{mathalfa}
\usepackage{xfrac}
\usepackage[unicode,hidelinks,bookmarks=false]{hyperref}
\newcommand{\ie}{i.e.\ }
\newcommand{\cf}{cf.\ }
\newcommand{\resp}{respectively\ }
\newcommand{\Subsection}{\subsection}
\providecommand{\nobreakdash}{\nobreak\hbox{-}}
\setlength{\emergencystretch}{2em}
\multlinegap0pt
\usepackage{amssymb}
\usepackage{float}
\usepackage{bm}
\usepackage[normalem]{ulem}
\newtheorem{theorem}{Theorem}
\newtheorem{lemma}{Lemma}
	\DeclareMathOperator{\Div}{div}
\def\Ga{\Gamma}
\def\Om{\Omega}
	\newcommand{\R}{\mathbb R}
\def\al{\alpha}
\def\be{\beta}
\def\bf{\textbf}
\def\de{\delta}
\def\df{\mathrm d}
\def\e{\varepsilon }
\def\f{\frac}
\def\ga{\gamma}
\def\mcH{\mathcal{H}}
\def\mcO{\mathcal{O}}
\def\ol{\overline}
\def\om{\omega}
\def\pt{\partial}
\def\s{\subseteq}
\def\se{\setminus}
	\DeclareMathOperator{\supp}{{supp}}
\def\ts{\times}
\def\vp{\varphi}
\title{\bf {\bf On Controllability of a Class of $N$-dimensional Hyperbolic Equations with Internal Single-point Degeneracy}}
\author{Donghui Yang, Weijia Wu}
\date{Source: arXiv:2605.05142v1 (CC BY 4.0)}
\begin{document}
\maketitle

\noindent\emph{Source and licence.} \bf {\bf On Controllability of a Class of $N$-dimensional Hyperbolic Equations with Internal Single-point Degeneracy}. Version arXiv:2605.05142v1 (CC BY 4.0), distributed by arXiv under the Creative Commons Attribution 4.0 International licence, http://creativecommons.org/licenses/by/4.0/, as recorded in the arXiv licence metadata for this exact version and shown on the abstract page https://arxiv.org/abs/2605.05142v1. Full licence notice: \texttt{licenses/arxiv-2605.05142-CC-BY-4.0.txt}.\par\smallskip

\noindent\textit{Editorial scope.} Counted unit: the article's theorem Carleman2 (source0.tex lines 570--580). The article's complete original proof of it is delivered with it (source0.tex lines 953--1478, one \texttt{proof} environment of 526 lines, which the article places away from the statement). No mathematics is written, repaired or re-derived: the statement and the proof are the article's own lines, in the article's own notation, and no retained line is substituted or re-flowed.  Retained uncounted as the source-local context the proof needs: lines 516--522; lines 533--535; lines 559--565; lines 800--806; lines 814--835.  Nothing was omitted from any retained span, so the package carries no omission record.  No retained line contains a citation, so the wrapper carries no bibliography and no bibliography entry.  No retained line contains a figure environment, an image-inclusion command or drawing code, so no image is delivered and the package declares no asset dependency.  Editorial formatting only, in addition: an AMS \texttt{amsart} preamble that restates the article's own declarations and macros used by the retained lines, namely the environments \texttt{theorem}, \texttt{lemma} on separate counters, together with the standard packages those lines need.  The retained environments print in this document as Theorem 1, Lemma 1 in order of appearance, whereas the article prints the same items with the numbering of its own full document.  The complete native source of the article as distributed in the arXiv e-print for this exact version is bundled unchanged as \texttt{sources/199927/source0.tex}.
		\begin{equation}\label{EQ-1}
			\begin{cases}
				\pt_t^2 u-\Div\left(|x|^\alpha \nabla u\right)=\chi_\omega f, & \mbox{in } Q, \\
				u=0, & \mbox{on } \Sigma,\\
				u(0)=u_0, \pt_tu(0)=u_1, &\mbox{in }\Om,
			\end{cases}
		\end{equation}

		\begin{equation}\label{2-6}
			\mcO(\Ga_+, 3\de)\s \om,
		\end{equation}

		\begin{equation}\label{EQ-2}
			\begin{cases}
				\pt_t^2z-\Div(|x|^\al \nabla z)=F, &\mbox{in } Q,\\
				z=0, &\mbox{on }\Sigma,\\
				z(0)=z_0, \pt_tz(0)=z_1, &\mbox{in }\Om,
			\end{cases}
		\end{equation}

		\begin{lemma}\label{L2.3}
			Let $\Om\s \R^N$ be a domain with smooth boundary. Then, there exists a $C^2$-vector field $V$ such that
			\begin{equation*}
				V(x)=\nu(x), \ x\in \pt\Om, \ \mbox{ and } |V(x)|\leq 1, \ x\in \ol{\Om},
			\end{equation*}
			where $\nu$   is the unit outward normal vector to $\pt\Om$.
		\end{lemma}

		\begin{theorem}\label{T2.2}
			Given $T>0$, suppose that
			\begin{equation*}
				\chi_\omega f\in L^2(Q), \quad u_0\in \mcH_0^1(\Om), \ \mbox{ and } u_1\in L^2(\Om).
			\end{equation*}
			Then, there exists a unique solution $u$ to  \eqref{EQ-1} such that
			\begin{equation*}
				u\in C([0,T]; \mcH_0^1(\Om))\cap C^1([0,T]; L^2(\Om)),
			\end{equation*}
			and there exists a constant $C>0$ satisfying
			\begin{equation}\label{10.02.1}
				\|u\|_{C([0,T]; \mcH_0^1(\Om))}+\|\pt_tu\|_{C([0,T]; L^2(\Om))}\leq C\left(\|u_0\|_{\mcH_0^1(\Om)}+\|u_1\|_{L^2(\Om)}+\|\chi_\omega f\|_{L^2(Q)}\right).
			\end{equation}
			Furthermore,
			\begin{equation*}
				\pt_\nu u\in L^2(\Sigma),
			\end{equation*}
			and there exists a constant $C=C(T, \Om)>0$ such that
			\begin{equation}\label{10.02.4}
				\|\pt_\nu u\|_{L^2(\Sigma)}\leq C\left(\|u_0\|_{\mcH_0^1(\Om)}+\|u_1\|_{L^2(\Om)}+\|\chi_\omega f\|_{L^2(Q)}\right).
			\end{equation}
		\end{theorem}

		\begin{theorem}\label{Carleman2}
			Let $\om\subset\Om$ be given as in \eqref{2-6}. Assume  $0\in\om$. For any solution $z$ to \eqref{EQ-2}, there exist positive constants  $C=C(\om,\Om)$, such that for sufficiently large constants
			$s$ and $\gamma$, the following inequality holds:
			\begin{equation}\label{2.3}
				\begin{split}
					&\iint_Qe^{2s\vp}s\ga \vp \left(|x|^\al \nabla z\cdot \nabla z+|z_t|^2+s^2\ga^2\vp^2|z|^2\right)\df x\df t\\
					&\leq C\iint_{\om\ts (0,T)}e^{2s\vp}s\ga \vp\left( z_t^2+|x|^\al \nabla z\cdot \nabla z +s^2\ga^2\vp^2|z|^2 \right)\df x\df t+C\iint_Qe^{2s\vp} F^2\df x\df t,
				\end{split}
			\end{equation}
	where     $\varphi$  is the weight function defined in \eqref{vp}.
		\end{theorem}

		\begin{proof}[{\rm \bf { Proof of  Theorem \ref{Carleman2}.}}]
			Since $0\in\om$, we assume there exists $\varepsilon >0$ such that  $B(0,3\varepsilon )\s\om$. Let
			\begin{equation*}
				z\in C^1([0,T]; L^2(\Om))\cap C([0,T]; \mcH_0^1(\Om))
			\end{equation*}
			be a solution of the following adjoint system (i.e., \eqref{EQ-2}):
			\begin{equation*}
				\begin{cases}
					\pt_t^2z-\Div\left(|x|^\al \nabla z\right)=F, & \mbox{in } Q,\\
					z(0)=z_0, \pt_tz(0)=z_1, &\mbox{in }\pt Q,\\
					z=0, &\mbox{on }\Sigma
				\end{cases}
			\end{equation*}
			with
			\begin{equation}\label{*}
				\supp z\s \Om\se B(0,\e).
			\end{equation}
			Let
			\begin{equation*}
				\psi(x,t)=|x|^2-\be(t-t_0)^2+\be_0
			\end{equation*}
			with $\be\in (0,\varrho), t_0\in (0,T),\be_0\geq 0$. Define
			\begin{equation}\label{vp}
				\vp(x,t)=e^{\ga\psi(x,t)},
			\end{equation}
			and denote
			\begin{equation*}
				w=e^{s\vp}z
			\end{equation*}
			with $z$ being  a solution of \eqref{EQ-2}. It is straightforward to show that
			\begin{equation*}
				e^{s\vp}z_t=w_t-s\vp_tw, \quad e^{s\vp}z_{tt}=w_{tt}-2s\vp_tw_t+s^2\vp_t^2w-s\vp_{tt}w,
			\end{equation*}
			\begin{equation*}
				e^{s\vp}\nabla z=\nabla w-sw\nabla\vp,
			\end{equation*}
			and
			\begin{equation*}
				e^{s\vp}\Div(|x|^\al \nabla z)=\Div\left(|x|^\al \nabla w\right)-2s|x|^\al \nabla \vp\cdot \nabla w-sw\Div\left(|x|^\al \nabla\vp\right)+s^2|x|^\al w\nabla\vp\cdot \nabla\vp.
			\end{equation*}
			These relations imply that
			\begin{equation*}
				\begin{split}
					e^{s\vp}\left(\pt_t^2z-\Div\left(|x|^\al \nabla z\right)\right)
					&=w_{tt}-\Div\left(|x|^\al \nabla w\right)\\
					&\hspace{4.5mm}+s^2\left(\vp_t^2w-|x|^\al w\nabla\vp\cdot \nabla\vp\right)\\
					&\hspace{4.5mm}+s\left(-2\vp_tw_t-\vp_{tt}w+2|x|^\al \nabla\vp\cdot \nabla w+w\Div(|x|^\al \nabla\vp)\right).
				\end{split}
			\end{equation*}
			Hence,
			\begin{equation*}
				P^+w+P^-w=e^{s\vp}F,
			\end{equation*}
			with
			\begin{equation*}
				P^+w:=\left(w_{tt}-\Div\left(|x|^\al \nabla w\right)\right)+s^2\left(\vp_t^2-|x|^\al \nabla\vp\cdot \nabla\vp\right)w,
			\end{equation*}
			and
			\begin{equation*}
				P^-w:=-2s\left(\vp_tw_t-|x|^\al \nabla\vp\cdot \nabla w\right)-s\left(\vp_{tt}-\Div(|x|^\al \nabla\vp)\right)w.
			\end{equation*}
			Therefore,
			\begin{equation*}
				\begin{split}
					\|e^{s\vp}F\|^2=\|P^+w\|^2+\|P^-w\|^2+2(P^+w, P^-w).
				\end{split}
			\end{equation*}
			
			
			Now, we proceed to compute the following expression
			\begin{eqnarray*}
				(P^+w, P^-w)
				&&=-2s\iint_Q w_{tt}(\vp_tw_t-|x|^\al \nabla\vp\cdot \nabla w)\df x\df t\\ %1
				&& -s\iint_Qw_{tt}(\vp_{tt}-\Div (|x|^\al \nabla\vp))w\df x\df t\\ %2
				&&+2s\iint_Q \Div(|x|^\al \nabla w) \left(\vp_tw_t-|x|^\al \nabla\vp\cdot \nabla w\right)\df x\df t\\ %3
				&& +s\iint_Q\Div(|x|^\al \nabla w)\left(\vp_{tt}-\Div (|x|^\al \nabla \vp)\right)w\df x\df t\\ %4
				&&-2s^3\iint_Q(\vp_t^2w-|x|^\al w\nabla\vp\cdot \nabla\vp)\left(\vp_t w_t-|x|^\al \nabla\vp\cdot \nabla w\right)\df x\df t\\ %5
				&& -s^3\iint_Q\left(\vp_t^2-|x|^\al \nabla\vp\cdot \nabla\vp\right) \left(\vp_{tt}-\Div (|x|^\al \nabla\vp)\right)w^2 \df x\df t\\ %6
				&&=:\sum_{i=1}^6 I_i.
			\end{eqnarray*}
			
			
			
			Next, we compute  $I_1$. Given that  $w_t(0)=w_t(T)=0$ and $w|_{\Sigma}=0$, we derive
			\begin{eqnarray*}
				I_1
				&&=-2s\iint_Q w_{tt}(\vp_tw_t-|x|^\al \nabla\vp\cdot \nabla w)\df x\df t \\
				&&=-s\iint_Q \vp_t \pt_t w_t^2\df x\df t+2s\iint_Q |x|^\al \nabla\vp\cdot \nabla w w_{tt}\df x\df t\\
				&&=-s\iint_Q\pt_t\left(\vp_t w_t^2\right)\df x\df t+s\iint_Q \vp_{tt}w_t^2\df x\df t\\
				&& +2s\iint_Q\pt_t\left(|x|^\al \nabla\vp\cdot \nabla w w_t\right)\df x\df t-2s\iint_Q|x|^\al\nabla\vp_t \cdot \nabla w w_t\df x\df t-2s\iint_Q |x|^\al\nabla \vp\cdot \nabla w_t w_t\df x\df t\\
				&&=s\iint_Q \vp_{tt}w_t^2\df x\df t-2s\iint_Q|x|^\al(\nabla\vp_t \cdot \nabla w) w_t\df x\df t\\
				&& -s\iint_Q \Div\left(|x|^\al\nabla \vp    w_t^2\right) \df x\df t+s\iint_Q \Div\left(|x|^\al \nabla \vp\right) w_t^2\df x\df t\\
				&&=s\iint_Q \vp_{tt}w_t^2\df x\df t-2s\iint_Q|x|^\al(\nabla\vp_t \cdot \nabla w) w_t\df x\df t+s\iint_Q \Div\left(|x|^\al \nabla \vp\right) w_t^2\df x\df t\\
				&&=s\iint_Q\left(\vp_{tt}+\Div \left(|x|^\al \nabla\vp\right)\right) w_t^2\df x\df t-2s\iint_Q|x|^\al(\nabla\vp_t \cdot \nabla w) w_t\df x\df t.
			\end{eqnarray*}
			Now, let's compute $I_2$.  Given that $w(0)=w(T)=0$, we obtain
			\begin{eqnarray*}
				I_2
				&&=-s\iint_Q(\vp_{tt}-\Div (|x|^\al \nabla\vp))w_{tt}w\df x\df t\\
				&&=-s\iint_Q\pt_t\Big( \left(\vp_{tt}-\Div (|x|^\al \nabla\vp)\right)w_{t}w\Big)\df x\df t\\
				&& +s\iint_Q\left(\vp_{tt}-\Div (|x|^\al \nabla \vp)\right)w_t^2\df x\df t+\f{s}{2}\iint_Q \pt_t w^2\left(\pt_t^2\vp_t-\Div(|x|^\al \nabla \vp_t)\right)\df x\df t\\
				&&=s\iint_Q\left(\vp_{tt}-\Div (|x|^\al \nabla \vp)\right)w_t^2\df x\df t\\
				&&\hspace{4.5mm}+\f{s}{2}\iint_Q\pt_t\Big(w^2(\pt_t^2\vp_t-\Div(|x|^\al \nabla \vp_t))\Big)\df x\df t-\f{s}{2}\iint_Q w^2\left(\pt_{tt}^2\vp-\Div(|x|^\al \nabla \vp_{tt})\right)\df x\df t\\
				&&=s\iint_Q\left(\vp_{tt}-\Div (|x|^\al \nabla \vp)\right)w_t^2\df x\df t-\f{s}{2}\iint_Q w^2\left(\pt_{tt}^2\vp-\Div(|x|^\al \nabla \vp_{tt})\right)\df x\df t,
			\end{eqnarray*}
			where we have used the fact that the integral of a total derivative over a time interval $[0,T]$ with zero initial and final conditions vanishes.
			
			
			Next, we proceed to compute  $I_3$. Given the condition  $w|_{\Sigma}=0$,  we derive the following:
			\begin{eqnarray*}
				I_3
				&&=2s\iint_Q \Div(|x|^\al \nabla w) \left(\vp_tw_t-|x|^\al \nabla\vp\cdot \nabla w\right)\df x\df t\\
				&&=2s\iint_Q\Div\Big(|x|^\al \nabla w(\vp_tw_t-|x|^\al \nabla\vp\cdot \nabla w)\Big)\df x\df t-2s\iint_Q|x|^\al \nabla w\cdot \nabla (\vp_tw_t-|x|^\al \nabla\vp\cdot \nabla w)\df x\df t\\
				&&=-2s\iint_\Sigma |x|^{2\al}(\nabla w\cdot \nu)(\nabla \vp\cdot \nabla w)\df S\df t\\
				&&\hspace{4.5mm}-2s\iint_Q|x|^\al (\nabla w\cdot \nabla \vp_t)w_t\df x\df t-2s\iint_Q|x|^\al \vp_t\nabla w\cdot \nabla w_t\df x\df t\\
				&&\hspace{4.5mm}+2s\iint_Q|x|^{2\al-2}(\nabla w\cdot x)(\nabla \vp\cdot \nabla w)\df x\df t+2s\iint_Q|x|^{2\al}\nabla w\cdot \nabla(\nabla \vp\cdot \nabla w)\df x\df t\\
				&&=-s\iint_\Sigma  |x|^{2\al}(\nabla \vp\cdot\nu) |\pt_\nu w|^2 \df S\df t+2s\iint_Q|x|^{2\al-2}(\nabla w\cdot x)(\nabla \vp\cdot \nabla w)\df x\df t\\
				&&\hspace{4.5mm}-2s\iint_Q|x|^\al (\nabla w\cdot \nabla \vp_t)w_t\df x\df t+s\iint_Q|x|^\al \vp_{tt}(\nabla w\cdot \nabla w)\df x\df t\\
				&&\hspace{4.5mm}+2s\iint_Q|x|^{2\al}D^2\vp\nabla w\cdot \nabla w\df x\df t-s\iint_Q\Div(|x|^{2\al}\nabla\vp)|\nabla w|^2\df x\df t
			\end{eqnarray*}
			where we have used the identities:
			\begin{equation*}
				\begin{split}
					&-2s\iint_Q|x|^\al \vp_t\nabla w\cdot \nabla\vp_t\df x\df t\\
					&=-s\iint_Q|x|^\al\vp_t\pt_t\left(\nabla w\cdot \nabla w\right)\df x\df t\\
					&=-s\iint_Q\pt_t\left(|x|^\al \vp_t (\nabla w\cdot \nabla w)\right)\df x\df t+s\iint_Q|x|^\al \vp_{tt}(\nabla w\cdot \nabla w)\df x\df t\\
					&=s\iint_Q|x|^\al \vp_{tt}(\nabla w\cdot \nabla w)\df x\df t,
				\end{split}
			\end{equation*}
			and
			\begin{eqnarray*}
				&&2s\iint_Q|x|^{2\al}\nabla w\cdot \nabla(\nabla \vp\cdot \nabla w)\df x\df t\\
				&&=2s\iint_Q|x|^{2\al}D^2\vp\nabla w\cdot \nabla w\df x\df t+s\iint_Q|x|^{2\al}\nabla \vp\cdot \nabla (|\nabla w|^2)\df x\df t\\
				&&=2s\iint_Q|x|^{2\al}D^2\vp\nabla w\cdot \nabla w\df x\df t\\
				&&\hspace{4.5mm}+s\iint_Q\Div\left( |x|^{2\al}\nabla \vp |\nabla w|^2\right)\df x\df t-s\iint_Q\Div(|x|^{2\al}\nabla\vp)|\nabla w|^2\df x\df t\\
				&&=2s\iint_Q|x|^{2\al}D^2\vp\nabla w\cdot \nabla w\df x\df t\\
				&&\hspace{4.5mm}+s\iint_\Sigma  |x|^{2\al}(\nabla \vp\cdot\nu) |\pt_\nu w|^2 \df S\df t-s\iint_Q\Div(|x|^{2\al}\nabla\vp)|\nabla w|^2\df x\df t\\
			\end{eqnarray*}
			with $D^2\vp=(\pt_{x_ix_j}^2\vp)_{i,j\in J_N}$ and
			\begin{equation*}
				\nabla w\cdot \nabla (\nabla w\cdot \nabla \vp)=D^2\vp\nabla w\cdot \nabla w+\f{1}{2}\nabla \vp\cdot \nabla (|\nabla z|^2).
			\end{equation*}
			
			
			
			Now, let us compute  $I_4$. Given that  $w|_{\Sigma}=0$, we obtain
			\begin{eqnarray*}
				I_4
				&&=s\iint_Q\Div(|x|^\al \nabla w)\left(\vp_{tt}-\Div (|x|^\al \nabla \vp)\right)w\df x\df t\\
				&&=s\iint_Q\Div\Big(|x|^\al \nabla w\left(\vp_{tt}-\Div (|x|^\al \nabla \vp)\right)w\Big)\df x\df t-s\iint_Q|x|^\al \nabla w\cdot \nabla \Big( \left(\vp_{tt}-\Div (|x|^\al \nabla \vp)\right)w\Big)\df x\df t\\
				&&=-s\iint_Q|x|^\al \nabla w\cdot \nabla w(\vp_{tt}-\Div(|x|^\al \nabla \vp))\df x\df t-\f{s}{2}\iint_Q|x|^\al  \nabla (\vp_{tt}-\Div (|x|^\al \nabla\vp))\cdot \nabla w^2\df x\df t\\
				&&=-s\iint_Q|x|^\al \nabla w\cdot \nabla w(\vp_{tt}-\Div(|x|^\al \nabla \vp))\df x\df t\\
				&&\hspace{4.5mm}-\f{s}{2}\iint_Q\Div\Big(  \nabla (\vp_{tt}-\Div (|x|^\al \nabla\vp))w^2\Big)\df x\df t+\f{s}{2}\iint_Q\Div\Big(|x|^\al \nabla  \left(\vp_{tt}-\Div (|x|^\al \nabla \vp)\right) \Big) w^2\df x\df t\\
				&&=-s\iint_Q|x|^\al \nabla w\cdot \nabla w(\vp_{tt}-\Div(|x|^\al \nabla \vp))\df x\df t+\f{s}{2}\iint_Q\Div\Big(|x|^\al \nabla  \left(\vp_{tt}-\Div (|x|^\al \nabla \vp)\right) \Big) w^2\df x\df t.
			\end{eqnarray*}
			
			
			
			Now, we proceed to compute $I_5$. Given  $w(0)=w(T)=0$ and $ w|_{\Sigma}=0$, we derive:
			\begin{eqnarray*}
				I_5
				&&=-2s^3\iint_Q(\vp_t^2w-|x|^\al w\nabla\vp\cdot \nabla\vp)\left(\vp_t w_t-|x|^\al \nabla\vp\cdot \nabla w\right)\df x\df t\\
				&&=-2s^3\iint_Q\vp_t^3w_tw\df x\df t+2s^3\iint_Q|x|^\al \vp_t^2\nabla \vp\cdot \nabla w w\df x\df t\\
				&&\hspace{4.5mm}+2s^3\iint_Q |x|^\al \vp_t(\nabla \vp\cdot \nabla \vp)w_t w\df x\df t-2s^3\iint_Q |x|^{2\al}(\nabla\vp\cdot \nabla\vp)(\nabla \vp\cdot \nabla w)w\df x\df t \\
				&&=-s^3\iint_Q\vp_t^3\pt_tw^2\df x\df t+s^3\iint_Q|x|^\al\vp_t^2\nabla \vp\cdot \nabla w^2\df x\df t\\
				&&\hspace{4.5mm}+s^3\iint_Q|x|^\al \vp_t(\nabla\vp\cdot \nabla\vp)\pt_tw^2\df x\df t-s^3\iint_Q|x|^{2\al}(\nabla\vp\cdot \nabla\vp)\nabla\vp\cdot \nabla w^2\df x\df t\\
				&&=-s^3\iint_Q\pt_t\left(\vp_t^3w^2\right)\df x\df t+3s^3\iint_Q \vp_t^2\vp_{tt}w^2\df x\df t\\
				&&\hspace{4.5mm}+s^3\iint_Q\Div\Big(|x|^\al \vp_t^2\nabla \vp w\Big)\df x\df t-s^3\iint_Q\Div(|x|^\al \vp_t^2\nabla \vp)w^2\df x\df t\\
				&&\hspace{4.5mm}+s^3\iint_Q\pt_t\Big(|x|^\al\vp_t(\nabla\vp\cdot \nabla \vp) w\Big)\df x\df t-s^3\iint_Q\pt_t(|x|^\al \vp_t(\nabla \vp\cdot \nabla \vp))w^2\df x\df t\\
				&&\hspace{4.5mm}-s^3\iint_Q\Div\Big(|x|^{2\al}(\nabla\vp\cdot \nabla\vp)\nabla\vp w^2\Big)\df x\df t+s^3\iint_Q\Div\left(|x|^{2\al}(\nabla\vp\cdot \nabla\vp)\nabla\vp\right)w^2\df x\df t\\
				&&7=3s^3\iint_Q \vp_t^2\vp_{tt}w^2\df x\df t-s^3\iint_Q\Div(|x|^\al \vp_t^2\nabla \vp)w^2\df x\df t\\
				&&\hspace{4.5mm}-s^3\iint_Q\pt_t(|x|^\al \vp_t(\nabla \vp\cdot \nabla \vp))w^2\df x\df t+s^3\iint_Q\Div\left(|x|^{2\al}(\nabla\vp\cdot \nabla\vp)\nabla\vp\right)w^2\df x\df t.
			\end{eqnarray*}
			
			Next, we compute $I_6$:
			\begin{equation*}
				\begin{split}
					I_6=&-s^3\iint_Q\left(\vp_t^2-|x|^\al \nabla\vp\cdot \nabla\vp\right) \left(\vp_{tt}-\Div (|x|^\al \nabla\vp)\right)w^2 \df x\df t.
				\end{split}
			\end{equation*}
			
			Combining  $I_1$ to $I_6$, we obtain
			\begin{eqnarray*}
				&&(P^+w,P^-w)\\
				&&=s\iint_Q\left(\vp_{tt}+\Div \left(|x|^\al \nabla\vp\right)\right) w_t^2\df x\df t-2s\iint_Q|x|^\al(\nabla\vp_t \cdot \nabla w) w_t\df x\df t\\ % I_1\\
				&&\hspace{4.5mm}+s\iint_Q\left(\vp_{tt}-\Div (|x|^\al \nabla \vp)\right)w_t^2\df x\df t-\f{s}{2}\iint_Q w^2\left(\pt_{tt}^2\vp -\Div(|x|^\al \nabla \vp_{tt})\right)\df x\df t\\ %I_2\\
				&&\hspace{4.5mm}-s\iint_\Sigma  |x|^{2\al}(\nabla \vp\cdot\nu) |\pt_\nu w|^2 \df S\df t+2s\iint_Q|x|^{2\al-2}(\nabla w\cdot x)(\nabla \vp\cdot \nabla w)\df x\df t\\
				&&\hspace{4.5mm}-2s\iint_Q|x|^\al (\nabla w\cdot \nabla \vp_t)w_t\df x\df t+s\iint_Q|x|^\al \vp_{tt}(\nabla w\cdot \nabla w)\df x\df t\\
				&&\hspace{4.5mm}+2s\iint_Q|x|^{2\al}D^2\vp\nabla w\cdot \nabla w\df x\df t-s\iint_Q\Div(|x|^{2\al}\nabla\vp)|\nabla w|^2\df x\df t\\ %I_3\\
				&&\hspace{4.5mm} -s\iint_Q|x|^\al \nabla w\cdot \nabla w(\vp_{tt}-\Div(|x|^\al \nabla \vp))\df x\df t+\f{s}{2}\iint_Q\Div\Big(|x|^\al \nabla  \left(\vp_{tt}-\Div (|x|^\al \nabla \vp)\right) \Big) w^2\df x\df t \\ %I_4
				&&\hspace{4.5mm} +3s^3\iint_Q \vp_t^2\vp_{tt}w^2\df x\df t-s^3\iint_Q\Div(|x|^\al \vp_t^2\nabla \vp)w^2\df x\df t\\
				&&\hspace{4.5mm}-s^3\iint_Q\pt_t(|x|^\al \vp_t(\nabla \vp\cdot \nabla \vp))w^2\df x\df t+s^3\iint_Q\Div\left(|x|^{2\al}(\nabla\vp\cdot \nabla\vp)\nabla\vp\right)w^2\df x\df t\\ % I_5
				&&\hspace{4.5mm} -s^3\iint_Q\left(\vp_t^2-|x|^\al \nabla\vp\cdot \nabla\vp\right) \left(\vp_{tt}-\Div (|x|^\al \nabla\vp)\right)w^2 \df x\df t\\
				&&=\sum_{i=1}^{17} H_i.
			\end{eqnarray*}
			
			Now, we estimate $(P^+w,P^-w)$. The fifth term $H_5$  does not need estimation:
			\begin{equation*}
				-s\iint_\Sigma  |x|^{2\al}(\nabla \vp\cdot\nu) |\pt_\nu w|^2 \df S\df t =-2s\ga \iint_\Sigma |x|^{2\al}  \vp (x\cdot \nu)|\pt_\nu w|^2\df S\df t.
			\end{equation*}
			
			Next, we estimate the sixth term $H_6$:
			\begin{eqnarray*}
				&&2s\iint_Q|x|^{2\al-2}(\nabla w\cdot x)(\nabla \vp\cdot \nabla w)\df x\df t\\
				&&=2s\iint_Q |x|^{2\al-2}(\nabla w\cdot x) (2\ga \vp x\cdot \nabla w)\df x\df t\\
				&&=4s\ga\iint_Q|x|^{2\al-2}\vp (x\cdot \nabla w)^2\df x\df t.
			\end{eqnarray*}
			
			Now, we estimate the tenth term $H_{9}$:
			\begin{eqnarray*}
				H_{9}
				&&=2s\iint_Q|x|^{2\al}D^2\vp\nabla w\cdot \nabla w\df x\df t\\
				&&=2s\iint_Q|x|^{2\al} \vp\left(2\ga \de_{ij}+4\ga^2x_ix_j\right) \nabla w\cdot \nabla w\df x\df t\\
				&&=2s \alpha \ga \iint_Q|x|^{2\al}\vp |\nabla w|^2\df x\df t+8s\ga^2\iint_Q|x|^{2\al}(x\cdot \nabla w)^2\df x\df t.
			\end{eqnarray*}
			
			
			Next, we estimate the first term $H_1$:
			\begin{eqnarray*}
				H_1
				&&=s\iint_Q\left(\vp_{tt}+\Div \left(|x|^\al \nabla\vp\right)\right) w_t^2\df x\df t\\
				&&=s\iint_Q\vp(-2\be\ga+\ga^2\psi_t^2)w_t^2\df x\df t+s\iint_Q\Div(|x|^\al \nabla \vp)w_t^2\df x\df t\\
				&&\geq-2s\be\ga\iint_Q\vp w_t^2\df x\df t+s\ga^2\iint_Q \psi_t^2\vp w_t^2\df x\df t+2\ga s\iint_Q\vp\left(N+1+2\ga |x|^2\right)|x|^\al w_t^2\df x\df t.
			\end{eqnarray*}
			
			
			Now, we estimate the second term $H_2$:
			\begin{eqnarray*}
				H_2
				&&=-2s\iint_Q|x|^\al(\nabla\vp_t \cdot \nabla w) w_t\df x\df t\\
				&&=-2s\ga^2\iint_Q |x|^\al \vp  \psi_t\nabla \psi \cdot \nabla w w_t\df x\df t=-2s\ga^2\iint_Q( |x|^\al \vp^\f{1}{2}x\cdot \nabla w)(\vp^\f{1}{2}\psi_tw_t)\df x\df t \\
				&&\geq -s\ga^2\iint_Q|x|^{2\al}\vp(x\cdot \nabla w)^2\df x\df t-s\ga^2 \iint_Q\psi_t^2 \vp w_t^2\df x\df t.
			\end{eqnarray*}
			
			
			Now, we estimate $H_3$:
			\begin{eqnarray*}
				H_3
				&& =s\iint_Q\left(\vp_{tt}-\Div (|x|^\al \nabla \vp)\right)w_t^2\df x\df t\\
				&&\geq -2s\be\ga \iint_Q\vp w_t^2\df x\df t+s\ga^2\iint_Q \psi_t^2\vp w_t^2\df x\df t-2\ga s\iint_Q\vp\left(N+1+2\ga |x|^2\right)|x|^\al w_t^2\df x\df t.
			\end{eqnarray*}
			
			
			Now, we estimate $H_7$:
			\begin{equation*}
				H_7
				=-2s\iint_Q|x|^\al (\nabla w\cdot \nabla \vp_t)w_t\df x\df t=H_2.
			\end{equation*}
			
			Next, we estimate  $H_8+H_{10}+H_{11}$:
			\begin{eqnarray*}
				&&H_8+H_{10}+H_{11}\\
				&&=s\iint_Q|x|^\al \vp_{tt}(\nabla w\cdot \nabla w)\df x\df t -s\iint_Q\Div(|x|^{2\al}\nabla\vp)|\nabla w|^2\df x\df t\\
				&&\hspace{4.5mm}-s\iint_Q|x|^\al \nabla w\cdot \nabla w(\vp_{tt}-\Div(|x|^\al \nabla \vp))\df x\df t\\
				&&=-s\iint_Q \nabla (|x|^\al) |x|^\al \nabla \vp|\nabla w|^2\df x\df t=-2s \alpha\ga\iint_Q|x|^{2\al}\vp|\nabla w|^2\df x\df t.
			\end{eqnarray*}
			
			
			Now, we obtain that
			\begin{equation*}
				\begin{split}
					& H_1+H_2+H_3+H_5+H_6+H_7+H_8+H_9+H_{10}+H_{11}\\
					&\geq 2(2-\alpha)s\ga \iint_Q|x|^{2\al}\vp |\nabla w|^2\df x\df t+6s\ga^2\iint_Q|x|^{2\al}(x\cdot \nabla w)^2\df x\df t\\
					&\hspace{4.5mm}+4s\ga\iint_Q|x|^{2\al-2}\vp (x\cdot \nabla w)^2\df x\df t \\
					&\hspace{4.5mm}-2s\ga \iint_\Sigma |x|^{2\al}  \vp (x\cdot \nu)|\pt_\nu w|^2\df S\df t -4s\be\ga\iint_Q\vp w_t^2\df x\df t.
				\end{split}
			\end{equation*}
			
			
			Now, we compute
			\begin{eqnarray*}
				&&(P^+w, s\ga \vp w)\\
				&& =\Big(\left(w_{tt}-\Div\left(|x|^\al \nabla w\right)\right)+s^2\left(\vp_t^2-|x|^\al \nabla\vp\cdot \nabla\vp\right)w, s\ga \vp w\Big)\\
				&&=s\ga \iint_Q \vp w_{tt} w\df x\df t-s\ga \iint_Q \vp \Div(|x|^\al \nabla w)w\df x\df t\\
				&&\hspace{4.5mm} +s^3\ga \iint_Q \vp_t^2\vp w^2\df x\df t-s^3\ga \iint_Q |x|^\al \vp (\nabla\vp\cdot \nabla\vp) w^2\df x\df t\\
				&&=\sum_{i=1}^4L_i.
			\end{eqnarray*}
			
			
			Now, we compute $L_1$. Since $w(0)=w(T)=0$, we obtain
			\begin{eqnarray*}
				L_1
				&&=s\ga \iint_Q \vp w_{tt} w\df x\df t\\
				&&=s\ga\iint_Q \pt_t(\vp w_tw)\df x\df t-s\ga\iint_Q \vp w_t^2\df x\df t-s\ga \iint_Q\vp_tw_tw\df x\df t\\
				&&=-s\ga\iint_Q\vp w_t^2\df x\df t -\f{s\ga }{2}\iint_Q\vp_t\pt_tw^2\df x\df t\\
				&&=-s\ga\iint_Q\vp w_t^2\df x\df t-\f{s\ga }{2}\iint_Q \pt_t(\vp_tw^2)\df x\df t+\f{s\ga }{2}\iint_Q \vp_{tt}w^2\df x\df t\\
				&&=-s\ga\iint_Q\vp w_t^2\df x\df t+\f{s\ga }{2}\iint_Q \vp_{tt}w^2\df x\df t.
			\end{eqnarray*}
			
			
			
			Now, we compute $L_2$:
			\begin{equation*}
				\begin{split}
					L_2
					&=-s\ga \iint_Q \vp \Div(|x|^\al \nabla w)w\df x\df t\\
					&=-s\ga \iint_Q\Div\left(|x|^\al \vp \nabla w w\right)\df x\df t+s\ga\iint_Q|x|^\al \nabla\vp\cdot \nabla w w\df x\df t+s\ga\iint_Q|x|^\al \vp \nabla w\cdot \nabla w\df x\df t\\
					&=\f{s\ga}{2}\iint_Q|x|^\al \nabla \vp\cdot \nabla w^2\df x\df t+s\ga \iint_Q|x|^\al \vp\nabla w\cdot \nabla w\df x\df t\\
					&=\f{s\ga}{2}\iint_Q\Div\left(|x|^\al \nabla \vp w^2\right)\df x\df t-\f{s\ga}{2}\iint_Q\Div\left(|x|^\al \nabla \vp\right)w^2\df x\df t+s\ga\iint_Q|x|^\al \vp \nabla w\cdot \nabla w\df x\df t\\
					&=-\f{s\ga}{2}\iint_Q\Div\left(|x|^\al \nabla \vp\right)w^2\df x\df t+s\ga\iint_Q|x|^\al \vp \nabla w\cdot \nabla w\df x\df t.
				\end{split}
			\end{equation*}
			
			Now, we compute $L_3+L_4$:
			\begin{eqnarray*}
				&&L_3+L_4\\
				&&=s^3\ga \iint_Q \vp_t^2\vp w^2\df x\df t-s^3\ga^3 \iint_Q |x|^\al \vp (\nabla\vp\cdot \nabla\vp) w^2\df x\df t=s^3\ga^3\iint_Q \vp^3(\psi_t^2-4|x|^{\al+2})w^2\df x\df t.
			\end{eqnarray*}
			Hence,
			\begin{eqnarray*}
				&&(P^+w, s\ga \vp w)\\
				&&=-s\ga\iint_Q\vp w_t^2\df x\df t+\f{s\ga }{2}\iint_Q \Big(\vp_{tt}-\Div\left(|x|^\al \nabla \vp\right)\Big)w^2\df x\df t\\
				&&\hspace{4.5mm} +s\ga\iint_Q|x|^\al \vp \nabla w\cdot \nabla w\df x\df t +s^3\ga^3\iint_Q \vp^3(\psi_t^2-4|x|^{\al+2})w^2\df x\df t,
			\end{eqnarray*}
			which implies that
			\begin{eqnarray*}
				s\ga \iint_Q\vp w_t^2\df x\df t
				&&\leq s^3\ga^3\iint_Q\vp^3\left|\psi_t^2-4|x|^{\al+2}\right| w^2\df x\df t+\f{s\ga}{\e^{\al}}\iint_Q |x|^{2\al}\vp\nabla w\cdot \nabla w\df x\df t\\
				&&\hspace{4.5mm}+\eta\|P^+w\|^2+\f{s^2\ga^3}{2\eta}\iint_Q  \vp^2w^2\df x\df t,
			\end{eqnarray*}
			i.e.,
			\begin{eqnarray*}
				&& s\ga \iint_Q|x|^{2\al}\vp|\nabla w|^2\df x\df t\\
				&&\geq s\ga \e^\al \iint_Q\vp w_t^2\df x\df t-\e^\al\left(s^3\ga^3\iint_Q\vp^3\left|\psi_t^2-4|x|^{\al+2}\right|w^2\df x\df t-\eta\|P^+z\|^2+\f{s^2\ga^2}{2\eta}\iint_Q\vp^2w^2\df x\df t\right).
			\end{eqnarray*}
			Then we get
			\begin{equation*}
				\begin{split}
					& H_1+H_2+H_3+H_5+H_6+H_7+H_8+H_9+H_{10}+H_{11}\\
					&\geq (2-\alpha)s\ga \iint_Q|x|^{2\al}\vp |\nabla w|^2\df x\df t+6s\ga^2\iint_Q|x|^{2\al}(x\cdot \nabla w)^2\df x\df t\\
					&\hspace{4.5mm}+4s\ga\iint_Q|x|^{2\al-2}\vp (x\cdot \nabla w)^2\df x\df t +s\ga \be\iint_Q\vp w_t^2\df x\df t\\
					&\hspace{4.5mm}-2s\ga \iint_\Sigma |x|^{2\al}  \vp (x\cdot \nu)|\pt_\nu w|^2\df S\df t\\
					&\hspace{4.5mm}-(2-\alpha)\e^\al\left(s^3\ga^3\iint_Q\vp^3\left|\psi_t^2-4|x|^{\al+2}\right|w^2\df x\df t-\eta\|P^+z\|^2+\f{s^2\ga^3}{2\eta}\iint_Q\vp^2w^2\df x\df t\right).
				\end{split}
			\end{equation*}
			by choosing  $0<\be<\f{2-\alpha}{5}\e^\al$.
			
			
			
			Now, we estimate $H_4+H_{12}$.  It is straightforward to derive that
			\begin{eqnarray*}
				&& -\f{s}{2}\iint_Q w^2\left(\pt_{tt}^2\vp -\Div(|x|^\al \nabla \vp_{tt})\right)\df x\df t+\f{s}{2}\iint_Q\Div\Big(|x|^\al \nabla  \left(\vp_{tt}-\Div (|x|^\al \nabla \vp)\right) \Big) w^2\df x\df t\\
				&&\geq -\f{s\ga^4}{2}\iint_Q(\psi_t^2-4|x|^{\al+2})^2\vp w^2\df x\df t-Cs\ga^3T^3\iint_Q \vp w^2\df x\df t.
			\end{eqnarray*}
			
			
			Next, we estimate  $H_{13}+H_{14}+H_{15}+H_{16}+H_{17}$:
			\begin{eqnarray*}
				&& H_{13}+H_{14}+H_{15}+H_{16}+H_{17}\\
				&& =+3s^3\iint_Q \vp_t^2\vp_{tt}w^2\df x\df t-s^3\iint_Q\Div(|x|^\al \vp_t^2\nabla \vp)w^2\df x\df t\\
				&&\hspace{4.5mm}-s^3\iint_Q\pt_t(|x|^\al \vp_t(\nabla \vp\cdot \nabla \vp))w^2\df x\df t+s^3\iint_Q\Div\left(|x|^{2\al}(\nabla\vp\cdot \nabla\vp)\nabla\vp\right)w^2\df x\df t\\
				&&\hspace{4.5mm} -s^3\iint_Q\left(\vp_t^2-|x|^\al \nabla\vp\cdot \nabla\vp\right) \left(\vp_{tt}-\Div (|x|^\al \nabla\vp)\right)w^2 \df x\df t\\
				&&=2s^3\iint_Q\vp_t^2\vp_{tt}w^2\df x\df t-4s^3\iint_Q|x|^\al \vp_t\nabla\vp_t\cdot \nabla \vp w^2\df x\df t\\
				&&\hspace{4.5mm}+s^3\iint_Q |x|^\al \nabla \vp\cdot \nabla \left(|x|^\al \nabla \vp\cdot \nabla \vp\right)w^2\df x\df t\\
				&&=2s^3\ga^3\iint_Q\vp^3\psi_t^2\psi_{tt}w^2\df x\df t+2s^3\ga^4\iint_Q\vp^3\psi_t^4w^2\df x\df t-16s^3\ga^4\iint_Q|x|^{\al+2}\psi_t^2\vp^3w^2\df x\df t\\
				&&\hspace{4.5mm}+8(\al+2)s^3\ga^3\iint_Q|x|^{2\al+2}\vp^3w^2\df x\df t+32s^3\ga^4\iint_Q|x|^{2\al+4}\vp^3w^2\df x\df t\\
				&&=-4s^3\ga^3\be\iint_Q (\psi_t^2-4|x|^{\al+2})\vp^3w^2\df x\df t+2s^3\ga^4\iint_Q\left(\psi_t^2-4|x|^{\al+2}\right)^2\vp^3w^2\df x\df t\\
				&&\hspace{4.5mm}+8(\al+2)s^3\ga^3\iint_Q|x|^{2\al+2}\vp^3w^2\df x\df t-16s^3\ga^3\be\iint_Q |x|^{\al+2}\vp^3w^2\df x\df t.
			\end{eqnarray*}
			
			
			
			Overall, we obtain
			\begin{eqnarray*}
				&&(P^+w,P^-w)\\
				&&\geq (2-\alpha)s\ga \iint_Q|x|^{2\al}\vp |\nabla w|^2\df x\df t+6s\ga^2\iint_Q|x|^{2\al}(x\cdot \nabla w)^2\df x\df t\\
				&&\hspace{4.5mm}+4s\ga\iint_Q|x|^{2\al-2}\vp (x\cdot \nabla w)^2\df x\df t +s\ga \be\iint_Q\vp w_t^2\df x\df t\\
				&&\hspace{4.5mm}-2s\ga \iint_\Sigma |x|^{2\al}  \vp (x\cdot \nu)|\pt_\nu w|^2\df S\df t\\
				&&\hspace{4.5mm}-(2-\alpha)\e^\al\left(s^3\ga^3\iint_Q\vp^3\left|\psi_t^2-4|x|^{\al+2}\right|w^2\df x\df t-\eta\|P^+z\|^2+\f{s^2\ga^3}{2\eta}\iint_Q\vp^2w^2\df x\df t\right)\\
				&&\hspace{4.5mm} -\f{s\ga^4}{2}\iint_Q(\psi_t^2-4|x|^{\al+2})^2\vp w^2\df x\df t-\f{Cs\ga^3T^3}{\e^{2\al+2}}\iint_Q |x|^{2\al+2} \vp w^2\df x\df t\\
				&&\hspace{4.5mm}-4s^3\ga^3\be\iint_Q \left|\psi_t^2-4|x|^{\al+2}\right|\vp^3w^2\df x\df t+2s^3\ga^4\iint_Q\left(\psi_t^2-4|x|^{\al+2}\right)^2\vp^3w^2\df x\df t\\
				&&\hspace{4.5mm}+8(\al+2)s^3\ga^3\iint_Q|x|^{2\al+2}\vp^3w^2\df x\df t-16s^3\ga^3\be\iint_Q |x|^{\al+2}\vp^3w^2\df x\df t.
				3	
			\end{eqnarray*}
			Now, choose $\eta=\f{1}{4\e^\al}$. Note that
			\begin{eqnarray*}
				&&2s^3\ga^4\iint_Q(\psi_t^2-4|x|^{\al+2})^2\vp^3w^2\df x\df t-(\e^\al +4\be)s^3\ga^3\iint_Q\vp^3\left|\psi_t^2-4|x|^{\al+2}\right| w^2\df x\df t\\
				&&=2s^3\ga^4\iint_Q(\psi_t^2-4|x|^{\al+2})^2\vp^3w^2\df x\df t- \f{(\e^\al+4\be)^2}{4\e^{2\al+2}}s^3\ga^3\iint_Q\left(\psi_t^2-4|x|^{\al+2}\right)^2\vp^3w^2\df x\df t\\
				&&\hspace{4.5mm} +\iint_Q\left[\f{\e^\al+4\be}{2\e^{\al+1}}\left|\psi_t^2-4|x|^{\al+2}\right|-\e^{\al+1}\right]^2\vp^3w^2\df x\df t-\e^{2\al+2}s^3\ga^3\iint_Q\vp^3w^2\df x\df t\\
				&&\geq 2s^3\ga^4\iint_Q(\psi_t^2-4|x|^{\al+2})^2\vp^3w^2\df x\df t- \f{(\e^\al+4\be)^2}{4\e^{2\al+2}}s^3\ga^3\iint_Q\left(\psi_t^2-4|x|^{\al+2}\right)^2\vp^3w^2\df x\df t\\
				&&\hspace{4.5mm} +\iint_Q\left[\f{\e^\al+4\be}{2\e^{\al+1}}\left|\psi_t^2-4|x|^{\al+2}\right|-\e^{\al+1}\right]^2\vp^3w^2\df x\df t- s^3\ga^3\iint_Q|x|^{2\al+2}\vp^3w^2\df x\df t.
			\end{eqnarray*}
			Select $\ga^*=\ga_*(\e)>1, s_*(\ga)>1$   sufficiently large, then for all $\ga \geq \ga_*, s\geq s_*(\ga)$, we obtain
			\begin{eqnarray*}
				&&(P^+w,P^-w)\\
				&&\geq (2-\alpha)s\ga \iint_Q|x|^{2\al}\vp |\nabla w|^2\df x\df t+6s\ga^2\iint_Q|x|^{2\al}(x\cdot \nabla w)^2\df x\df t\\
				&&\hspace{4.5mm}+4s\ga\iint_Q|x|^{2\al-2}\vp (x\cdot \nabla w)^2\df x\df t +s\ga \be\iint_Q\vp w_t^2\df x\df t \\
				&&\hspace{4.5mm}-2s\ga \iint_\Sigma |x|^{2\al}  \vp (x\cdot \nu)|\pt_\nu w|^2\df S\df t-\f{1}{4}\|P^+z\|^2\\
				&&\hspace{4.5mm}+s^3\ga^4\iint_Q\left(\psi_t^2-4|x|^{\al+2}\right)^2\vp^3w^2\df x\df t +4(\al+2)s^3\ga^3\iint_Q|x|^{2\al+2}\vp^3w^2\df x\df t.
			\end{eqnarray*}
			These imply that
			\begin{eqnarray*}
				&&2(2-\alpha)s\ga \iint_Q|x|^{2\al}\vp |\nabla w|^2\df x\df t+12s\ga^2\iint_Q|x|^{2\al}(x\cdot \nabla w)^2\df x\df t\\
				&&+8s\ga\iint_Q|x|^{2\al-2}\vp (x\cdot \nabla w)^2\df x\df t +2s\ga \be\iint_Q\vp w_t^2\df x\df t \\
				&&+2s^3\ga^4\iint_Q\left(\psi_t^2-4|x|^{\al+2}\right)^2\vp^3w^2\df x\df t +8(\al+2)s^3\ga^3\iint_Q|x|^{2\al+2}\vp^3w^2\df x\df t\\
				&&\leq 4s\ga \iint_{\Sigma} |x|^{2\al}  \vp (x\cdot \nu)|\pt_\nu w|^2\df S\df t +\|e^{s\vp}F\|^2\\
				&&\leq 4s\ga \iint_{\Sigma_+} |x|^{2\al}  \vp (x\cdot \nu)|\pt_\nu w|^2\df S\df t +\|e^{s\vp}F\|^2.
			\end{eqnarray*}
			Hence, we have the estimate
			\begin{equation}\label{A-2}
				\begin{split}
					&\iint_Qs\ga \vp \left(|x|^\al \nabla w\cdot \nabla w+|w_t|^2+s^2\ga^2\vp^2|w|^2\right)\df x\df t\\
					&\leq C \iint_Qe^{2s\vp}F^2\df x\df t+C\iint_{\Sigma_+} s\ga \vp|\pt_\nu w|^2\df S\df t.
				\end{split}
			\end{equation}
			
			
			Now, we  estimate
			\begin{equation*}
				\iint_{\Sigma_+}  s\ga \vp  |\pt_\nu w|^2\df S\df t\leq C  \int_0^T\int_{\mcO(\Ga_+,3\de)}   s\ga \vp\left(|x|^\al \nabla u\cdot \nabla u+w_t^2+w^2\right)\df x\df t .
			\end{equation*}
			
			
			
			By Lemma \ref{L2.3}, there exists a $C^2$-vector field $V$ such that
			\begin{equation*}
				V(x)=\nu(x), \ x\in\pt\Om, \ \mbox{ and } |V(x)|\leq 1, \ x\in\Om.
			\end{equation*}
			Choose a function $\rho(x)\in C^2(\ol\Om)$ satisfying
			\begin{equation*}
				0\leq \rho\leq 1 \mbox{ on }\R^N,\quad \rho\equiv 1 \mbox{ on } \mcO(\Ga_+, \de)\cap \Om, \quad \rho=0 \mbox{ on }\Om\se \mcO(\Ga_+, 2\de).
			\end{equation*}
			Set
			\begin{equation*}
				g=V\rho .
			\end{equation*}
			Multiplying the first equality of \eqref{EQ-2} by  $g\cdot \nabla w$ and integrating over  $Q$, we obtain
			\begin{equation*}
				\begin{split}
					\iint_Q\left(\pt_t^2w-\Div(|x|^\al \nabla w)\right) (g\cdot \nabla w)\df x\df t=\iint_Q F(g\cdot \nabla w)\df x\df t.
				\end{split}
			\end{equation*}
			
			
			The proof follows a similar approach to the proof of Theorem \ref{T2.2}. Since
			$w(0)=w(T)=0$ and $w|_{\Sigma}=0$,  we derive
			\begin{eqnarray*}
				\iint_Q \pt_t^2 w(g\cdot \nabla w)\df x\df t
				&&=\iint_Q\pt_t\left(w_t(g\cdot \nabla w)\right)\df x\df t-\f{1}{2}\iint_Qg\cdot \nabla w_t^2\df x\df t \\
				&&=-\f{1}{2}\iint_Q \Div(g w_t^2)\df x\df t+\f{1}{2}\iint_Q (\Div g) w_t^2\df x\df t =\f{1}{2}\iint_Q (\Div g) w_t^2\df x\df t.
			\end{eqnarray*}
			Since
			\begin{eqnarray*}
				&&-\iint_Q\Div(|x|^\al \nabla w)(g\cdot \nabla w)\df x\df t\\
				&&=-\iint_Q\Div\left(|x|^\al \nabla w (g\cdot \nabla w)\right)\df x\df t+\iint_Q|x|^\al \nabla w\cdot \nabla (g\cdot \nabla w)\df x\df t\\
				&&=-\iint_\Sigma |x|^\al (\nabla w\cdot \nu)(g\cdot \nabla w)\df S\df t+\iint_Q |x|^\al \left((Dg\nabla w)\cdot \nabla w+\f{1}{2} g\cdot \nabla |\nabla w|^2\right)\df x\df t\\
				&&=-\iint_{\Sigma} |x|^\al \vp |\pt_\nu w|^2\df x\df t+ \iint_Q |x|^\al (Dg\nabla w)\cdot \nabla w\df x\df t\\
				&&\hspace{4.5mm}+\f{1}{2}\iint_Q\Div\left(|x|^\al g|\nabla w|^2\right)\df x\df t-\f{1}{2}\iint_Q \Div(|x|^\al g)|\nabla w|^2\df x\df t\\
				&&=-\f{1}{2}\iint_{\Sigma} |x|^\al \vp |\pt_\nu w|^2\df x\df t+ \iint_Q |x|^\al (Dg\nabla w)\cdot \nabla w\df x\df t-\f{1}{2}\iint_Q \Div(|x|^\al g)|\nabla w|^2\df x\df t.
			\end{eqnarray*}
			From the above, using the definition of $\rho$, we obtain
			\begin{equation}\label{10.03.7}
				\begin{split}
					\iint_{\Sigma_+} \vp|\pt_\nu w|^2\df x\df t
					&\leq C\iint_{\mcO(\Ga_+,2\de)\ts (0,T)} \vp\left(w_t^2+|x|^\al \nabla w\cdot \nabla w\right)\df x\df t+C\iint_Q\vp F^2\df x\df t\\
					&\leq C\iint_{\om\ts (0,T)}\vp\left(w_t^2+|x|^\al \nabla w\cdot \nabla w\right)\df x\df t+C\iint_Q\vp F^2\df x\df t.
				\end{split}
			\end{equation}
			Combining this with \eqref{A-2} and using classical arguments, we can transform back to the original variable $z$ and conclude the following result:
			\begin{equation}\label{4.5}
				\begin{split}
					&\iint_Qe^{2s\vp}s\ga \vp \left(|x|^\al \nabla z\cdot \nabla z+|z_t|^2+s^2\ga^2\vp^2|z|^2\right)\df x\df t\\
					&\leq C\iint_{\om\ts (0,T)}e^{2s\vp}s\ga \vp\left( z_t^2+|x|^\al \nabla z\cdot \nabla z +s^2\ga^2\vp^2|z|^2 \right)\df x\df t+C\iint_Qe^{2s\vp} F^2\df x\df t.
				\end{split}
			\end{equation}
			
			
			It is worth noting that the aforementioned conclusions are contingent upon the establishment of assumption \ref{*}. Subsequently, we shall delve into the discussion of acquiring equation \eqref{2.3} in the absence of assuming Condition \ref{*}.
			
			Let$z=z\kappa + z(1-\kappa)$, $\hat{z}=z(1-\kappa)$ , where
			$$
			\kappa \in C^\infty (\Omega), \quad \kappa \equiv 1 \mbox{ in } B\left(0,\frac{\varepsilon }{2}\right), \quad \kappa \equiv 0 \mbox{ in } \Om\backslash B(0,\varepsilon ), \quad 0\le \kappa \le 1.
			$$
			Then   $\hat{z} \subset \Om\backslash B(0,\varepsilon )$, and
			\begin{eqnarray*}
				\pt_t^2 \hat{z}-\Div\left(|x|^\alpha \nabla \hat{z}\right)
				&&=  \pt_t^2 z(1-\kappa)-\Div\left(|x|^\alpha \nabla\left[  z(1-\kappa)\right] \right)\\
				&&= \pt_t^2 z(1-\kappa)-\Div\left(|x|^\alpha \left[ \nabla z(1-\kappa) + z \nabla (1-\kappa) \right] \right)\\
				&&= \pt_t^2 z(1-\kappa)-\Div\left(|x|^\alpha \nabla z(1-\kappa)\right) -\Div\left(|x|^\alpha z \nabla (1-\kappa)  \right)\\
				&&= \pt_t^2 z(1-\kappa)-\Div\left(|x|^\alpha \nabla z \right)(1-\kappa) + |x|^\alpha \nabla z \cdot \nabla \kappa + \Div\left(|x|^\alpha z \nabla \kappa  \right)\\
				&&= F(1-\kappa) + |x|^\alpha \nabla z \cdot \nabla \kappa + \Div\left(|x|^\alpha z \nabla \kappa  \right)\\
				&&=:\hat{F}.
			\end{eqnarray*}
			Now, given that we have $\pt_t^2 \hat{z}-\Div\left(|x|^\alpha \nabla \hat{z}\right)=\hat{F}$ and $\hat{z}$
			satisfies \eqref{*}, we derive the following inequality:
			\begin{equation}\label{4.6}
				\begin{split}
					&\iint_Qe^{2s\vp}s\ga \vp \left(|x|^\al \nabla \hat{z}\cdot \nabla \hat{z}+|\hat{z}_t|^2+s^2\ga^2\vp^2|\hat{z}|^2\right)\df x\df t\\
					&\leq C\iint_{\om\ts (0,T)}e^{2s\vp}s\ga \vp\left( \hat{z}_t^2+|x|^\al \nabla \hat{z}\cdot \nabla \hat{z} +s^2\ga^2\vp^2|\hat{z}|^2 \right)\df x\df t+C\iint_Qe^{2s\vp} \hat{F}^2\df x\df t.
				\end{split}
			\end{equation}
			From this, we can deduce that
			\begin{eqnarray*}
				\iint_Qe^{2s\vp}s\ga \vp |x|^\al \nabla z\cdot \nabla z dx dt
				&=&\iint_Qe^{2s\vp}s\ga \vp |x|^\al \nabla \left( z\xi + \hat{z}\right) \cdot \nabla \left( z\xi + \hat{z}\right) dx dt\\
				&\le& \iint_Qe^{2s\vp}s\ga \vp |x|^\al \nabla \hat{z}\cdot \nabla \hat{z} dx dt + \iint_Qe^{2s\vp}s\ga \vp |x|^\al \nabla \left(z\xi \right) \cdot \nabla \left(z\xi \right) dx dt\\
				&+& 2 \iint_Qe^{2s\vp}s\ga \vp |x|^\al \nabla \left(z\xi \right) \cdot \nabla \hat{z} dx dt\\
				&\le& 2\iint_Qe^{2s\vp}s\ga \vp |x|^\al \nabla \hat{z}\cdot \nabla \hat{z} dx dt + 2\iint_Qe^{2s\vp}s\ga \vp |x|^\al \nabla \left(z\xi \right) \cdot \nabla \left(z\xi \right) dx dt.
			\end{eqnarray*}
			Similarly, we obtain
			\begin{equation*}
				\begin{split}
					\iint_Qe^{2s\vp}s\ga \vp |z_t|^2 dx dt
					&\le 2\iint_Qe^{2s\vp}s\ga \vp |(z \xi)_t | ^2 dx dt + 2\iint_Qe^{2s\vp}s\ga \vp |\hat{z}_t|^2  dx dt,
				\end{split}
			\end{equation*}
			and
			\begin{equation*}
				\begin{split}
					\iint_Qe^{2s\vp}s\ga \vp |z|^2 dx dt
					&\le 2\iint_Qe^{2s\vp}s\ga \vp |z\xi|^2 dx dt + 2\iint_Qe^{2s\vp}s\ga \vp |\hat{z}| ^2 dx dt,
				\end{split}
			\end{equation*}
			Utilizing \eqref{4.6} and the properties of  $\kappa$, we derive
			\begin{eqnarray*}
				&&\iint_Qe^{2s\vp}s\ga \vp \left(|x|^\al \nabla z\cdot \nabla z+|z_t|^2+s^2\ga^2\vp^2|z|^2\right)\df x\df t\\
				&&\le C\iint_Qe^{2s\vp}s\ga \vp \left(|x|^\al \nabla \hat{z}\cdot \nabla \hat{z}+|\hat{z}_t|^2+s^2\ga^2\vp^2|\hat{z}|^2\right)\df x\df t\\
				&&+
				C\iint_Qe^{2s\vp}s\ga \vp \left(|x|^\al \nabla (z \xi)\cdot \nabla (z \xi)+|(z \xi)_t|^2+s^2\ga^2\vp^2|z \xi|^2\right)\df x\df t\\
				&&\le
				C\iint_{\om\ts (0,T)}e^{2s\vp}s\ga \vp\left( \hat{z}_t^2+|x|^\al \nabla \hat{z}\cdot \nabla \hat{z} +s^2\ga^2\vp^2|\hat{z}|^2 \right)\df x\df t+C\iint_Qe^{2s\vp} \hat{F}^2\df x\df t\\
				&&+C\iint_{\om\ts (0,T)}e^{2s\vp}s\ga \vp\left( z_t^2+|x|^\al \nabla z\cdot \nabla z +s^2\ga^2\vp^2|z|^2 \right)\df x\df t\\
				&&\le C\iint_{\om\ts (0,T)}e^{2s\vp}s\ga \vp\left( z_t^2+|x|^\al \nabla z\cdot \nabla z +s^2\ga^2\vp^2|z|^2 \right)\df x\df t +C\iint_Qe^{2s\vp} \hat{F}^2\df x\df t.
			\end{eqnarray*}
			From this, we can deduce \eqref{2.3}. 		
		\end{proof}

\end{document}
